\chapter{Group B Streptococcus}

\textbf{Current guideline note:} In the United States, routine prenatal screening is performed at 36 0/7--37 6/7 weeks of gestation. This current timing supersedes the older 35--37-week wording retained in some historical references.

\section*{What Is This Test?}
Group B Streptococcus (GBS) testing detects \textit{Streptococcus agalactiae}, a gram-positive, catalase-negative, beta-hemolytic coccus classified by Lancefield group B antigen. GBS commonly colonizes the gastrointestinal and genital tracts; approximately 10--30\% of pregnant women are colonized at a given time. Colonization is usually asymptomatic, but transmission during labor or after rupture of membranes can cause early-onset neonatal disease, including bacteremia, pneumonia, and meningitis. GBS can also cause late-onset disease and invasive infection in adults, particularly older adults and people with diabetes or immunocompromising conditions. Universal prenatal screening and intrapartum antibiotic prophylaxis reduce, but do not eliminate, early-onset neonatal GBS disease.

\section*{Alternative Names}
GBS Culture; Group B Strep Screen; GBS Screening; Vaginal-Rectal GBS Culture; GBS PCR; Intrapartum GBS NAAT; \textit{Streptococcus agalactiae} Culture; GBS Latex Agglutination

\section*{Principle}
Prenatal screening culture: Vaginal and rectal swab collected at 35--37 weeks' gestation. The swab is placed in a selective enrichment broth containing gentamicin and nalidixic acid (Lim broth or Strep B Carrot Broth) that inhibits gram-negative organisms and most other gram-positive organisms while promoting GBS growth. After 18--24 hours of incubation at 35--37\textdegree{}C, the broth is subcultured onto 5\% sheep blood agar and incubated for an additional 18--24 hours. Colonies of GBS produce a narrow zone of beta-hemolysis (some strains are non-hemolytic, ~5\%). Identification: CAMP test positive (arrowhead-shaped enhanced hemolysis when streaked perpendicular to \textit{Staphylococcus aureus}), hippurate hydrolysis, Lancefield grouping by latex agglutination (group B antigen), MALDI-TOF, or automated identification systems. Pigment production in Granada medium (orange/red colonies) is specific for GBS. Chromogenic agars (StrepB Select, CHROMagar StrepB) produce pink/mauve colonies. Sensitivity of enrichment culture (the gold standard) is >95\%. Intrapartum NAAT: Cepheid GeneXpert GBS or Roche cobas GBS performs real-time PCR targeting the \textit{cfb} gene (CAMP factor) directly from a vaginal-rectal swab without enrichment, providing results in 1--2 hours. Sensitivity 95--98\%, specificity 95--99\%. Intrapartum NAAT is indicated for women with unknown GBS status presenting in labor (no prenatal culture, or result unavailable). Rapid antigen tests for GBS from vaginal specimens have low sensitivity (40--70\%) and are NOT recommended for screening.

\section*{History}
GBS was first recognized as a cause of bovine mastitis in the 1880s. Human GBS infection was first described by Rebecca Lancefield in the 1930s after the development of the Lancefield grouping system. GBS emerged as the leading cause of neonatal sepsis in the 1970s. The landmark study by Carol Baker and colleagues in the 1970s--1980s established the relationship between maternal GBS colonization and neonatal sepsis. The first CDC guidelines for GBS prevention (1996) recommended either universal culture-based screening at 35--37 weeks or a risk-factor-based approach. The 2002 CDC guidelines adopted universal culture-based screening as the sole recommended strategy, refined in 2010 and 2019 (ACOG Committee Opinion 797). Intrapartum NAAT was introduced in the 2000s and recommended in the 2010 guidelines for women with unknown GBS status in labor. The American College of Obstetricians and Gynecologists (ACOG) and the CDC reaffirm universal screening and IAP as the standard of care.

\section*{Why This Test Is Ordered}
\begin{itemize}
  \item Universal prenatal screening of all pregnant women at 35--37 weeks' gestation (vaginal and rectal culture for GBS). This is the standard of care per ACOG and CDC. Screening at this gestational age maximizes the predictive value for GBS colonization at delivery. Colonization status can change, and cultures collected more than 5 weeks before delivery have reduced predictive accuracy
  \item Intrapartum GBS NAAT for women presenting in labor at term with unknown GBS status (no prenatal culture performed, culture results unavailable, or clinical scenario not covered by prenatal screening)
  \item Evaluation of women with preterm labor (<37 weeks) and unknown GBS status: perform GBS culture or NAAT on admission and administer IAP until results are available; if the culture is negative at 48 hours, IAP can be discontinued
  \item Diagnosis of neonatal GBS sepsis: Blood cultures, CSF culture (for meningitis), and CSF PCR from symptomatic neonates within the first 72 hours of life (early-onset) or between 7--90 days (late-onset)
  \item Diagnosis of invasive GBS disease in nonpregnant adults: Blood cultures, CSF cultures, and cultures from other normally sterile sites (synovial fluid, pleural fluid)
  \item Susceptibility testing for GBS isolates from penicillin-allergic patients: Clindamycin and erythromycin susceptibility is essential for guiding IAP in penicillin-allergic women at high risk for anaphylaxis (GBS resistance to clindamycin approaches 40--50\%). Vancomycin is the alternative when clindamycin resistance is demonstrated or unknown
\end{itemize}

\section*{Is This a Primary or Secondary Test}
GBS culture at 36 0/7--37 6/7 weeks is the primary screening test for most pregnant women. Intrapartum NAAT is an alternative for selected women with unknown GBS status in labor. Blood and CSF cultures are primary tests for suspected neonatal GBS sepsis. Rapid antigen tests are not recommended as a substitute for prenatal culture.

\section*{How This Test Is Performed}
Prenatal screening: A single swab (Dacron or rayon, not cotton) is used to sample first the lower vagina (vaginal introitus), then the rectum (insert through the anal sphincter). The patient should be instructed on self-collection or the healthcare provider collects the specimen. The swab is placed in a non-nutritive transport medium (Amies without charcoal is preferred) or directly into Lim broth enrichment medium and transported at room temperature. The specimen should be clearly labeled as a prenatal GBS screen to ensure the laboratory uses selective enrichment broth. For intrapartum NAAT, the same dual-site swab is collected and placed in the manufacturer-specific transport tube (for GeneXpert GBS, the swab is placed in the elution reagent vial). Blood cultures for neonatal GBS: 1--2 mL of blood inoculated into a pediatric blood culture bottle. CSF culture: 0.5--1 mL CSF in a sterile tube.

\section*{What Preparation Is Needed}
No special preparation for vaginal-rectal swab. The patient should be informed that both a vaginal and rectal sample will be taken (or self-collected). There is no discomfort from the rectal swab beyond mild sensation.

\section*{Contraindications and Precautions}
GBS culture at 35--37 weeks is the standard of care and should be performed on ALL pregnant women regardless of planned mode of delivery. Even women with planned cesarean delivery should be screened because: (a) cesarean delivery does not eliminate the risk of neonatal GBS infection (ascending infection may occur with rupture of membranes before cesarean); (b) GBS colonization increases the risk of maternal intrapartum/postpartum infection (chorioamnionitis, endometritis). Women with GBS bacteriuria during the current pregnancy (any colony count, not just >10$^5$ CFU/mL) should NOT be screened at 35--37 weeks---they are considered GBS-positive and automatically qualify for IAP (GBS bacteriuria indicates heavy colonization). Women with a previous infant with GBS disease automatically qualify for IAP in all subsequent pregnancies regardless of screening result. IAP is NOT indicated for cesarean delivery performed before the onset of labor and with intact membranes, regardless of GBS colonization status. Penicillin G is the preferred agent for IAP (5 million units IV loading dose, then 2.5--3 million units IV every 4 hours until delivery); ampicillin is an acceptable alternative. For penicillin-allergic women at low risk for anaphylaxis (non-IgE-mediated reaction), cefazolin is recommended. For penicillin-allergic women at high risk for anaphylaxis (IgE-mediated: urticaria, angioedema, anaphylaxis), clindamycin is the recommended agent ONLY if the GBS isolate is susceptible to both clindamycin and erythromycin. If resistance or unknown susceptibility, vancomycin is recommended.

\section*{Risks}
Vaginal-rectal swab collection is painless and poses no risk to the mother or fetus. The primary risk is NOT from testing but from failure to provide IAP to a GBS-colonized woman, which has an attack rate (risk of EOS in the neonate born to a GBS-colonized mother without IAP) of 1--2\%, with 5--10\% mortality and 20--30\% permanent neurologic sequelae among survivors of neonatal GBS meningitis.

\section*{What Happens After the Test}
Prenatal culture results available in 48--72 hours. Intrapartum NAAT results available in 1--2 hours. A positive result at 35--37 weeks: The patient is designated GBS-positive, and IAP is indicated at the time of labor or rupture of membranes. The result should be clearly documented in the medical record and communicated to the labor and delivery unit. A negative result at 35--37 weeks: No IAP indicated unless intrapartum risk factors develop (gestational age <37 weeks, rupture of membranes $\geq$18 hours, intrapartum fever $\geq$100.4\textdegree{}F/38\textdegree{}C, or prior infant with GBS disease). Women with unknown GBS status at the time of labor ($\geq$37 weeks) should receive IAP based on risk factors if NAAT is not available: gestational age <37 weeks, duration of membrane rupture $\geq$18 hours, or intrapartum temperature $\geq$100.4\textdegree{}F. For preterm labor (<37 weeks) with unknown GBS status: GBS culture or NAAT should be obtained and IAP initiated; it can be discontinued after 48 hours if cultures are negative.

\section*{Understanding Results}

\subsection*{Below Normal}
\begin{itemize}
  \item \textbf{Prenatal GBS culture negative (no GBS isolated):} GBS not detected in the vaginal or rectal specimen. No IAP is indicated at the time of delivery (assuming delivery at term, without prolonged rupture of membranes or intrapartum fever). Sensitivity is >95\% but not 100\%; a small risk of false-negative culture exists, but routine IAP for culture-negative women is not recommended
  \item \textbf{Intrapartum NAAT negative:} GBS DNA not detected. In a woman with unknown GBS status presenting in labor, a negative NAAT rules out GBS colonization and IAP is NOT indicated (assuming term delivery, no prolonged ROM, no intrapartum fever)
  \item \textbf{Neonatal blood and CSF cultures negative at 48--72 hours:} Makes neonatal GBS sepsis unlikely but does not exclude it definitively (antibiotics given to the mother during labor may suppress neonatal culture growth; low-level bacteremia may be missed with small blood volume). If clinical suspicion persists, continue empiric antibiotics
\end{itemize}

\subsection*{Above Normal}
\begin{itemize}
  \item \textbf{Prenatal GBS culture positive:} GBS colonization confirmed. The patient is a GBS carrier and should receive IAP at the time of labor (penicillin G preferred, or ampicillin). The positive result should be flagged prominently in the medical record and communicated to the labor unit. Oral antibiotics before labor to eradicate GBS colonization are NOT effective (GBS recolonizes rapidly) and should NOT be prescribed
  \item \textbf{Intrapartum NAAT positive:} GBS DNA detected. In a woman with unknown GBS status, IAP is indicated. The intrapartum NAAT obviates the need for empiric IAP in women with risk factors but no documented GBS colonization (provides a more rational approach than risk-factor-based prophylaxis)
  \item \textbf{GBS bacteriuria (any colony count) during pregnancy:} Indicates heavy GBS colonization. The patient is considered GBS positive for the remainder of the pregnancy and does NOT require a 35--37-week rectovaginal screen. IAP is indicated at delivery. Additionally, GBS bacteriuria (>10$^5$ CFU/mL) should be treated as a symptomatic UTI during pregnancy (penicillin or amoxicillin)
  \item \textbf{Neonatal blood culture positive for GBS:} GBS sepsis confirmed. Full evaluation for meningitis (lumbar puncture with CSF culture and PCR) is required even if the blood culture is positive, because meningitis coexists with bacteremia in 10--30\% of cases and requires a longer treatment course (14--21 days for meningitis vs. 10 days for bacteremia without meningitis)
  \item \textbf{GBS from normally sterile site in adult (blood, CSF, synovial fluid):} Invasive GBS disease. All GBS isolates from invasive disease should undergo antimicrobial susceptibility testing (penicillin MIC, clindamycin, erythromycin). Penicillin resistance in GBS is extremely rare, but when present (penicillin MIC >0.12 mcg/mL), it should be confirmed by a reference laboratory
\end{itemize}

\section*{Normal Results}
\textbf{GBS not isolated} from prenatal vaginal-rectal culture. \textbf{GBS DNA not detected} by intrapartum NAAT. Blood/CSF cultures: \textbf{No growth at 5 days}.

\section*{Follow-up Tests}
\begin{itemize}
  \item \textbf{GBS antimicrobial susceptibility testing (clindamycin and erythromycin):} On all GBS isolates from penicillin-allergic women at high risk for anaphylaxis, to determine whether clindamycin can be used for IAP. Clindamycin resistance (constitutive and inducible) approaches 40--50\%; erythromycin resistance predicts inducible clindamycin resistance (D-test positive)
  \item \textbf{Neonatal CSF culture and CSF PCR:} For all neonates with suspected GBS sepsis; lumbar puncture is indiciated for all infants with positive blood cultures because meningitis may be present without clinical signs
  \item \textbf{Neonatal complete blood count with differential, CRP, procalcitonin:} Adjunctive laboratory markers for neonatal sepsis evaluation. A normal CBC and CRP at 6--12 hours of age have a high negative predictive value (>99\%) for ruling out EOS in well-appearing infants whose mothers received adequate IAP. The Kaiser Permanente neonatal sepsis risk calculator integrates these parameters with clinical examination
  \item \textbf{GBS PCR from sterile sites:} When culture is negative but clinical suspicion remains high (e.g., partially treated meningitis), 16S rRNA PCR and sequencing or GBS-specific PCR on CSF may provide diagnostic confirmation
\end{itemize}

\section*{References}
\begin{enumerate}
  \item ACOG Committee Opinion No. 797: Prevention of group B streptococcal early-onset disease in newborns. Obstet Gynecol. 2020;135(2):e51-e72.
  \item Verani JR, McGee L, Schrag SJ. Prevention of perinatal group B streptococcal disease: revised guidelines from CDC, 2010. MMWR Recomm Rep. 2010;59(RR-10):1-36.
  \item Schuchat A. Group B streptococcus. Lancet. 1999;353(9146):51-56.
  \item Puopolo KM, Lynfield R, Cummings JJ; AAP Committee on Fetus and Newborn. Management of infants at risk for group B streptococcal disease. Pediatrics. 2019;144(2):e20191881.
  \item Stoll BJ, Hansen NI, Sánchez PJ, et al. Early onset neonatal sepsis: the burden of group B streptococcal and E. coli disease continues. Pediatrics. 2011;127(5):817-826.
\end{enumerate}

\clearpage
